%%%%%%%%%%%%%%%%%%%%%%%%%%%%%%%%%
% Discussion                    %
%%%%%%%%%%%%%%%%%%%%%%%%%%%%%%%%%

\section{Discussion}

This chapter interprets the Chapter 5 results in their academic, practical, and regulatory context: the sector diagnostic and maturity profile, the theoretical and managerial implications of the framework and the DataPilot tool, its position against the existing state of practice, the rigor of the pre-development work, and the threats to validity with their mitigations.

\subsection{Interpretation of the Key Findings}

The central finding is that the data maturity of the Tunisian banking sector is concentrated at the emerging level, the second of five. The sector has moved beyond the ad hoc stage but not yet reached the formalization, standardization, and repeatability that define the third level. The worked Global Maturity Index (GMI) of 2.19 out of 5, equivalent to 44\%, is representative: an institution at this level typically possesses isolated good practices but lacks the documented, approved, and consistently applied processes that would make them durable and auditable.

The diagnostic explains why the sector sits here rather than higher. Governance, which carries the largest weight at 25\%, is also the weakest: only 30\% of respondents report a formally approved data strategy, so roughly seven in ten work without a board-level data mandate. Because governance is the foundation on which the other dimensions rest, this deficit propagates across ownership, quality controls, and investment in architecture and skills. The data quality finding reinforces this: that a majority verify data before use looks like diligence but reads more accurately as distrust, since professionals verify because they cannot rely on the data, signalling the absence of automated upstream controls. The architecture findings, slow access and only partial centralization, show that even where governance intent exists the technical substrate does not yet support it. And the analytics finding quantifies the business cost: data-driven decision making remains a minority practice against a majority still relying on intuition, so the sector is not yet extracting decision value from its data at scale.

Together these findings describe a coherent, self reinforcing pattern rather than isolated weaknesses, which implies that piecemeal remediation in isolation is unlikely to raise overall maturity. The framework's weighted, multidimensional structure is a direct response: it makes the interdependence explicit and discourages the common error of investing in analytics tools while leaving governance and quality unaddressed. The diagnostic does not describe a failing sector but one that has begun the right activities without yet institutionalizing them, exactly the transition from the second to the third maturity level that the framework is designed to measure and guide.

The weights are themselves an interpretive statement about this interdependence. Governance receives the largest share, 25\%, because it is causally upstream: a strategy, an ownership model, and a compliance mandate determine whether quality controls are funded, architecture consolidated, and analytics investment directed at genuine business questions. Data quality, architecture, and analytics each carry 20\%, the operational layers through which governance intent is realized and decision value produced. Skills and culture carries the smallest, 15\%, not because it is unimportant but because it is assessed only through proxy indicators and its effect is mediated by the other four rather than expressed directly in auditable artifacts. The weighting thus encodes a theory of how maturity is built, and the profile, in which the heavily weighted governance dimension is also the weakest, explains why the index settles at the emerging level.

\subsection{The Perception versus Reality Gap}

A finding that emerged during the design of the scoring methodology, with implications beyond this project, is the perception versus reality gap: institutions consistently rate their own maturity above what their documentary evidence supports. A bank may declare it has a data strategy when what exists is an informal, undated set of intentions never approved at board level. The evidence requirement and the evidence cap neutralize this bias: any indicator scored above two out of five without a verifiable artifact is automatically capped at two. This turns the assessment from a declarative exercise into an evidence based one, so the GMI measures not how mature an institution believes it is but how mature it can prove it is, precisely the standard of proof a regulatory audit expects.

The gap is not merely a nuisance; it is diagnostically informative. The distance between the score an institution would assign itself and the score its evidence supports directly measures the formalization deficit separating the emerging level from the defined level: where the two coincide, practice is documented; where they diverge, a capability exists but has not been made durable, transferable, and auditable. Because the cap applies deterministically at each individual indicator rather than as a global adjustment, the assessment localizes this deficit precisely, showing exactly which controls, owners, or policies lack the artifact that would make them defensible. The tool thus turns a well known self-assessment bias into a targeted list of documentation gaps, each mapping to a concrete remediation action in the roadmap.

This also aligns the instrument with regulatory supervision: a BCT inspection requires not a declaration that a control exists but the artifact, the dated policy, the approval record, the lineage documentation, or the audit trail, and by making evidence the precondition for any score above two the framework meets that same evidentiary standard, so a favorable DataPilot result reflects readiness for the scrutiny an inspection involves. The tool does not, however, substitute for supervisory judgment: only the regulator can pronounce on compliance, and the instrument makes no claim to predict or replace that judgment.

\subsection{Theoretical and Academic Contributions}

The study makes three contributions to knowledge on data maturity assessment. The first is the demonstration that no single established framework suffices for a regulated, sector specific context, and that a structured, criteria based selection and hybridization of multiple frameworks produces a more complete instrument than any of its sources. The study screened fifteen candidate frameworks against five selection criteria and explicitly mapped what was retained and excluded from each of the five selected frameworks, providing a replicable method for constructing hybrid maturity models in other regulated sectors.

The second contribution is the formalization of an auditable scoring methodology. Much existing maturity literature stops at descriptive level definitions and does not specify how indicator responses aggregate into a defensible global score, nor how to prevent inflation. The methodology defines a transparent aggregation chain, the sub-dimension average, dimension average, and weighted global sum, and attaches explicit business rules: the evidence cap, the skip logic with a twenty percent ceiling, and the non-skippable status of regulatory indicators. It thereby closes the gap between a qualitative model and a quantitative, traceable measurement instrument.

The third contribution is the explicit alignment of a data maturity framework with a specific national regulatory text, the BCT Circulaire N\textdegree{}2025-08, combined with the international BCBS 239 standard. The thirteen regulatory indicators embedded in the framework, and their non-skippable status, make regulatory compliance a measurable component of maturity rather than a separate exercise. This integration of a maturity model with the author's mapping to a national regulation together with the internationally recognized BCBS 239 principles is, to the knowledge of this study, not present in the publicly available frameworks reviewed. The distribution of the thirteen indicators is itself meaningful: seven fall within governance, three within data quality, two within architecture, one within analytics, and none within skills and culture. This concentration in governance and quality expresses, at the level of the instrument, the regulator's own priority ordering, in which accountability, traceability, and the integrity of reported data precede tooling and culture. The framework thus does not treat compliance as an external checklist bolted onto a maturity model; it embeds the requirements at the exact dimensions where they belong, so a compliance failure and a maturity weakness point to the same underlying gap.

A further theoretical implication follows from these three contributions together. The prevailing literature separates three concerns this work unifies: the descriptive question of what capabilities define maturity, the measurement question of how to score them defensibly, and the regulatory question of whether the institution meets its obligations. By binding a hybridized capability model to a deterministic, auditable scoring method and to a specific regulatory mapping, the framework shows these can be addressed within a single coherent instrument rather than three disjoint exercises. This unification is the study's principal theoretical claim, and it is what allows a single number, the GMI, to carry a diagnostic, a measurement, and a compliance meaning at once.

\subsection{Practical and Managerial Implications}

For bank management, the framework and the DataPilot tool convert an abstract, frequently avoided question, how mature is our data capability, into a concrete, quantified, trackable figure. The implications are threefold. First, the GMI gives the board a single number to track across successive assessments, turning data maturity into a managed objective rather than an aspiration; because the index is deterministic and evidence based, a change between assessments reflects a real change in documented capability rather than a shift in respondent optimism, making it a defensible key performance indicator for a data transformation program. Second, the gap analysis and three phase improvement roadmap translate the diagnosis into a prioritized action plan. Each sub-dimension receives a priority from critical to low according to its score, then is sequenced by time horizon: a first phase, over three months, covers sub-dimensions below the emerging threshold or carrying an unmet regulatory obligation; a second covers those not yet at the defined level; a third, once a sub-dimension reaches that level, is devoted to optimization. Recommendations are concrete and dimension specific, so a governance weakness surfaces actions such as approving a board-level data strategy and appointing an accountable executive, while the same dimension once stronger surfaces optimization such as an annual KPI-linked strategy review. The sequence runs from the critical governance and compliance gaps first to the optimization that follows, reflecting the interdependence discussed above: remediation that begins with the upstream governance layer is the only ordering that lets the downstream investments in quality, architecture, and analytics hold. Third, the compliance view and exportable evidence package give management a defensible artifact to present to the regulator, cutting the cost and uncertainty of preparing for a BCT inspection and turning a reactive scramble into a standing, continuously maintained record.

For EY Consulting Tunisia, the framework and the tool constitute a reusable consulting asset. A standardized, evidence based instrument lets the firm deliver maturity diagnostics with consistent rigor across engagements, benchmark institutions against one another on a common scale, and ground its recommendations in a transparent, auditable methodology. It reduces the manual effort of each assessment and improves the comparability of results. Just as importantly, it de-risks the consulting relationship: because the scoring rules are explicit and the evidence cap prevents inflation, the firm can defend each score to a client and to a regulator on the same objective basis, protecting the credibility of the recommendations that follow. Reusability also amortizes the intellectual effort of framework design across many engagements rather than repeating it for each, so consultants concentrate their expertise on interpretation and remediation rather than on reconstructing an assessment method every time.

Delivering the tool as a multi-tenant platform, in the fourth iteration, extends these implications from the instrument to the way it is operated, turning it from a deliverable handed over at the end of an engagement into a service EY operates across engagements. EY maintains the master template centrally, so every engagement starts from the same calibrated instrument. Each bank receives its own copy, which its Admin may tailor to local vocabulary and emphasis without any ability to alter the scoring rules, the evidence cap, the skip ceilings, or the compliance thresholds; because every copy starts from the canonical master and EY reviews submissions across banks, local tailoring stays visible and results stay methodologically comparable. Group assessments distribute the questionnaire to the departments that actually hold the evidence, each dimension assigned to the department best placed to answer it, replacing a single proxy respondent with accountable contributions whose authorship is recorded, which a coordinator consolidates into one auditable submission per bank. These submissions accumulate over time into the comparable evidence base that the sector level benchmarking discussed below presupposes. What protects each bank in this shared setting is no longer local custody of the data but tenant isolation enforced in the database itself, confining every row to its own bank, completed by invite-only access and the immutability of submitted results.

For the regulator and the sector as a whole, a common assessment standard makes sector level benchmarking possible. If multiple institutions are assessed with the same framework and the same scoring rules, their results become comparable, which enables the identification of systemic weaknesses, such as the governance deficit the diagnostic revealed, and supports the design of sector wide remediation policy. Because thirteen of the indicators are drawn directly from the regulatory texts, an aggregate view across institutions would also give the supervisor a sector level reading of compliance with the circular, distinguishing requirements that are broadly met from those, such as the formal approval of a data strategy, where the whole sector lags. A shared, auditable instrument thus offers the regulator a more precise instrument of supervision than institution by institution declarations, and offers the sector a common language in which to compare its progress.

\subsection{Comparison with the Existing State of Practice}

The contribution is best understood by comparison with the frameworks from which it draws. Each excels on a single axis: DAMA-DMBOK on structural breadth, DCAM on regulatory alignment for finance, CMMI-DMM on the maturity scale, TDWI on evaluation mechanics, and Gartner on executive communication. None delivers all of these simultaneously, and none is operationalized as a usable tool aligned with the Tunisian regulatory environment. The framework does not replace these references; it composes them. Its novelty lies in the integration: a structurally complete model that scores in an auditable way, communicates to executives, and maps to a national regulation, delivered through a governed platform a bank operates with its own accounts and an advisor administers across institutions. This addresses the three research gaps identified in Chapter 2, each left open by the existing state of practice: the absence of a sector specific and regulation aligned framework, of an auditable scoring method, and of a usable tool.

This composition, rather than replacement, defines both the value and the honest limits of the contribution: each source was kept for the axis on which it is strongest and set aside where another was stronger, so the hybrid inherits all five strengths at once. What it adds beyond that inheritance is the connective tissue none of the sources provides, the deterministic aggregation chain, the evidence cap, and the author's regulatory mapping described among the contributions above. The result is not a larger checklist but a different kind of instrument, at once descriptive, measurable, auditable, and regulatorily grounded.

Against the state of practice in Tunisian banking specifically, the comparison is sharper. Without a shared instrument, institutions have assessed their data capability, when at all, through informal, non comparable self appraisals, and preparation for regulatory scrutiny has been a reactive, engagement specific effort. The framework replaces this with a single, reproducible, evidence based standard whose results are comparable across institutions and across time. This is the practical sense in which the contribution advances the state of practice: not by proposing a novel definition of maturity, which the literature already supplies in abundance, but by making maturity in this regulated context something that can be measured the same way twice, defended to a supervisor, and acted upon through a prioritized roadmap.

\subsection{Rigor of the Pre-Development Work}

The credibility of the artifacts rests largely on work completed before a single line of the tool was written. The conceptual phases, the benchmark, the framework design, and the scoring methodology, meet the standard of rigor expected of a research contribution and a professional consulting engagement, not a thin preamble to software development.

The first element is the breadth of the benchmark. Rather than adopting the first plausible reference model, the study screened fifteen international frameworks against five selection criteria: international credibility, banking relevance, scoring logic, complementarity, and BCT and BCBS 239 alignment. Each was scored out of fifteen in a full comparison matrix, every acceptance and rejection justified with a documented reason, and only five retained. The ten rejections are explicit, argued decisions, on grounds ranging from vendor lock-in and proprietary non-referenceable content to a scope too narrow to cover the five dimensions. The selection was not a mechanical threshold: a proprietary model with a competitive score was still rejected because it could not be referenced academically, showing the rule was principled rather than arithmetic.

The second element is that the framework is a genuine synthesis rather than an adoption. Each of the five selected frameworks was assigned a defined role in the hybrid, and for each the deliberately excluded parts are named with the reason. The model was engineered from complementary components, not borrowed from a single source and relabeled, and the traceability from every structural choice back to its source and exclusion decisions is preserved in the benchmark documentation.

The third element is that the design is evidence driven. Each of the four directly observable dimensions is justified by a specific finding of the diagnostic survey of 44 banking professionals; the fifth is justified by the survey's inability to observe it directly, the very finding that motivated its proxy-only treatment and lowest weight of 15\%. The structure and its 25/20/20/20/15 weighting are therefore anchored in field evidence from the sector rather than in intuition or in a transposition of foreign priorities: the weakest observed capability, governance, receives the heaviest weight, and the only indirectly observable dimension, skills and culture, the lightest. The weighting is thus a reasoned, evidence-informed choice rather than an arbitrary one; as the construct-validity discussion below acknowledges, it was not tested against alternative weightings and should be read as defensible rather than provably optimal.

The fourth element is the depth of the scoring rubrics. The framework does not stop at naming 47 indicators; each carries five explicit level descriptions, amounting to 235 distinct evaluation criteria that were individually written and calibrated. This is the step most maturity models omit, and it is precisely what converts a maturity concept into an auditable instrument: an assessor confronted with an artifact matches it against a written criterion rather than interpreting a vague label.

The fifth element is the anti-gaming and regulatory rigor built into the method itself: the evidence cap, the proxy-only rule for the fifth dimension, the twenty percent skip ceiling, and the non-skippable status of the thirteen regulatory indicators, whose mechanics are examined in the thesis passage and threats analysis below. What matters here is that these rules were designed during the conceptual phase, before implementation, as properties of the method rather than afterthoughts of the code.

The final element is method discipline. The conceptual phases were run as a structured consulting engagement: a four-step iterative method leading from benchmark through hybrid design and evidence-based scoring to tool development, a documented work breakdown covering the six phases from framing to closure, a stakeholder mapping formalized in an interest and influence matrix, and a formal SWOT analysis of the benchmark that explicitly weighed adopting an existing framework against designing a bespoke one before the design decision was taken. Taken together, these six elements show that the pre-development work was not preparatory in the weak sense but constitutive: the defensibility of the DataPilot scores discussed throughout this chapter is inherited directly from the rigor of the phases that preceded the tool.

If a single thesis runs through this chapter, it is that the instrument is most defensible precisely where it is constrained. Each property defended above is produced not by an added capability but by a deliberate restriction. The evidence cap bounds optimism: no claim can raise a score above two unless a documented artifact supports it. The skip ceiling bounds omission: at most two indicators may be set aside in Governance, two in Data Quality, and one in each of the remaining three dimensions, so a score cannot be reshaped by discarding inconvenient questions. The non-skippable status of the thirteen regulatory indicators bounds evasion: the regulatory core can never be removed from the computation. The explicit proxy label and the fifteen percent weight bound the least observable dimension, so the one part of the construct that cannot be verified directly has the least influence on the index. And renormalizing the weights over the answered dimensions bounds the effect of an unanswered dimension, so an incomplete assessment yields a proportionate figure rather than a distorted one. In each case the constraint is what converts a self assessment into a measurement: the score is trustworthy exactly where the method refuses to let it float. The threats examined below are correspondingly concentrated where no such constraint can reach, in the honesty and quality of the evidence itself. The same logic carries into the fourth iteration, where the platform is defensible for the same reason the score is. The role hierarchy, under which every invitation creates a role exactly one rank below the inviter, bank-scoped tenant isolation, and the immutability of submissions are restrictions of the same kind, together converting a shared platform into a defensible system of record. The single-user milestone's confidentiality argument, keeping every response in the assessor's own browser, is thereby superseded for everything except in-progress solo answers: what now protects a bank is not where its data sits but the isolation the database itself enforces on every row, itself a constraint rather than a capability.

\subsection{Threats to Validity}

Validity is examined along the four standard validity dimensions together with a researcher-bias threat, with the mitigations described for each.

\subsubsection{Construct Validity}

Construct validity concerns whether the framework measures what it claims, namely data maturity. The principal threat is that the choice of dimensions, sub-dimensions, and indicators could fail to capture the full construct. This was mitigated in three ways: the structure was derived from five established frameworks rather than a single source or intuition; the regulatory indicators were anchored directly in the text of the BCT Circulaire and BCBS 239; and each indicator was given an explicit five level rubric that defines precisely what is measured. The decomposition into five dimensions, twelve sub-dimensions, and 47 indicators is itself a construct validity mechanism, breaking a broad and contested concept into narrow, separately observable questions each with its own rubric, so no single ambiguous item stands in for an entire capability. The residual threat concerns the fifth dimension, skills and culture, which cannot be observed directly and is assessed through proxy indicators such as training budget, headcount ratio, and certifications; these are objective and auditable but measure investment in data culture rather than the culture itself, an acknowledged approximation. It was deliberately bounded in two ways: the dimension was assigned the smallest weight, 15\%, so the approximation has the least influence on the global index, and it was labeled explicitly as proxy based, so its interpretation is never mistaken for a direct measurement of culture. A related concession concerns the weighting itself. The 25\%, 20\%, 20\%, 20\%, and 15\% allocation is justified by the diagnostic evidence and by the regulatory centrality of governance, but it was not comparatively tested against alternative weightings. It should therefore be regarded as defensible rather than provably optimal, and the pilot assessments described in the future work are the natural setting for its calibration.

\subsubsection{Internal Validity}

Internal validity concerns whether the conclusions follow correctly from the data. The main threat is score inflation through declarative self assessment, the perception versus reality gap discussed above, mitigated by the evidence requirement and the evidence cap, which prevent any unverified claim from contributing a score above two. A second threat is aggregation consistency: an error in propagation from indicator to global score would invalidate the result. This was mitigated by validating the aggregation chain, confirming that a change at the indicator level propagates correctly through its sub-dimension and dimension to the GMI and that the dimension weights sum to one hundred percent. A third threat, specific to missing responses, is that skipping indicators could distort a dimension score by silently changing its denominator. This was mitigated by two rules: skipped and unanswered indicators are excluded from the averages rather than scored as zero, so an omission neither inflates nor deflates the result, and skips are capped at twenty percent of the indicators in each dimension. Renormalizing the global index over the weight of the answered dimensions applies the same principle at the top level, keeping a partial assessment proportionate rather than biased by a missing dimension. This renormalization is intended for a genuinely partial assessment, and it exposes one control the delivered platform does not yet enforce: because an entirely unanswered dimension drops out of the weighted sum, a group coordinator who left the weakest dimension unmapped could raise the reported index by omission rather than improvement. A finalized group submission should therefore be required to map all five dimensions, so renormalization keeps a partial solo assessment meaningful without becoming a route to omit a dimension from an official one. Enforcing full dimension coverage at finalization is a hardening step for future work. The non-skippable status of the thirteen regulatory indicators further protects internal validity, because the items most consequential for the conclusions can never be removed from the computation.

\subsubsection{External Validity}

External validity concerns the generalizability of the findings. The diagnostic is based on a survey of 44 banking professionals, which provides indicative rather than statistically exhaustive coverage of the sector, so the headline percentages should be read as directional evidence rather than precise population parameters. Two distinct claims to generalizability must be separated, since they carry different risk. The first, the generalizability of the sector diagnostic, is genuinely limited by the sample and is presented as directional. The second, the generalizability of the framework as a measurement instrument, does not depend on the survey at all: the framework is applied institution by institution, and its validity for a given bank rests on the evidence collected for that bank rather than on the survey sample. Confusing the two would overstate the risk, since a limited sample constrains what can be said about the sector but says nothing about whether the instrument measures a single institution correctly. The framework is specific to the banking sector and the Tunisian regulatory environment by design; transferring it to insurance or another jurisdiction would require re-anchoring the regulatory indicators to the relevant texts. This bounded scope is deliberate, not an oversight: the value of embedding thirteen indicators from a specific national circular is inseparable from their being specific. A version for another sector or jurisdiction would substitute the corresponding regulatory mapping while retaining the same aggregation and evidence logic, which are domain independent.

\subsubsection{Conclusion Validity}

Conclusion validity concerns whether the relationship between the assessment and the maturity it reports is reliable and reproducible. Because the scoring methodology is fully deterministic, two assessors who enter the same raw scores and the same evidence will, by construction, arrive at the same GMI. This computational reproducibility lets a change in the index over time be attributed to a real change in capability rather than a change of assessor or tool, making the figure usable as a tracked management indicator. It should not, however, be mistaken for inter-rater reliability in the stronger sense: two assessors may still read the same rubric and judge the same artifact differently when assigning the raw one-to-five score, so identical outputs are guaranteed only for identical inputs. Whether independent assessors would produce the same inputs for the same bank was not measured in this work and is identified as future work. The per-indicator rubric narrows this latitude at the point of scoring but remains a mitigation rather than a proof of reliability. The principal residual threat is the quality and honesty of the evidence supplied, which the tool cannot verify autonomously and which remains the assessor's responsibility. This threat is contained rather than eliminated: because every score above two must be justified by a named artifact retained in the exportable evidence package, a subsequent review or regulatory inspection can re-examine the same artifacts and confirm or overturn the score, placing the honesty of the evidence under external check even though the tool cannot adjudicate it. The platform delivered in the fourth iteration strengthens this further: every answer in a group assessment records who supplied it, and every finalized result is frozen as an immutable submission, so neither the provenance nor the content of a reported figure can drift after the fact.

The word auditable deserves scrutiny, since it is the headline adjective of the methodological contribution. Auditability is established here by construction rather than by a field audit: the aggregation chain is deterministic, every score above two is tied to a named artifact, and the complete set of responses, evidence references, and results is exportable as a single package an inspector could re-examine. Conducted on simulated data, the evaluation demonstrates that an inspection could retrace every figure in the index back to its evidence, not that an inspection has done so. The claim is a property of the design, strongly evidenced at that level, with confirmation under a real supervisory review deliberately deferred to the pilot assessments identified in the future work. Stating the boundary this precisely keeps the adjective honest: auditable here means built so that an audit would succeed, not already audited.

\subsubsection{Researcher Bias}

A threat that cuts across the four dimensions above is that a single author designed the framework, built the tool, seeded the demonstration data, and evaluated the results. This concentration of roles is a genuine researcher-bias threat: no independent party challenged the construct, scored a bank against the rubric, or judged whether an artifact was sufficient. Three factors contain, without eliminating, this threat. First, the instrument is deliberately auditable: an independent reviewer can retrace every displayed figure from the stored answers through the published rules and confirm or overturn it, placing the results under external check even though the evaluation was not independent. Second, the limitations of the evaluation, its reliance on seeded demonstration data and the absence of a field pilot or an automated test suite, are stated openly rather than hidden. Third, an independent evaluation, in which practitioners other than the author score real institutions and an inter-rater reliability study is carried out, is identified as the primary direction for future validation. Until then, the results should be read as a demonstration of a faithful and internally consistent instrument rather than an externally validated measurement.

\subsection{Conclusion}

This chapter interpreted the maturity deficit of the Tunisian banking sector as a coherent, self reinforcing pattern rooted in governance, articulated the theoretical and managerial value of the framework and the DataPilot tool, positioned the contribution relative to the existing frameworks, made explicit the rigor of the pre-development work from which the artifacts derive their defensibility, and examined the threats to validity with their mitigations. The next chapter concludes the report and outlines directions for future work.
